% Summarization cascade. The terminal tier is what makes the cascade total.
\begin{tikzpicture}[
    font=\footnotesize,
    node distance=5mm,
    tier/.style  = {draw, rounded corners=2pt, minimum height=7mm,
                    minimum width=32mm, align=center, line width=0.5pt},
    llm/.style   = {tier, fill=blue!7,  draw=blue!55},
    det/.style   = {tier, fill=green!8, draw=green!50!black},
    gate/.style  = {draw, diamond, aspect=2.4, inner sep=0.4mm, fill=black!3,
                    align=center, font=\scriptsize, line width=0.5pt},
    arr/.style   = {-{Latex[length=1.8mm,width=1.4mm]}, line width=0.6pt},
    note/.style  = {font=\scriptsize\itshape, text=black!55, align=left},
    fail/.style  = {note, text=red!55!black},
  ]

  \node[tier, fill=black!3] (req) {summarization request};
  \node[llm, below=of req]  (ollama) {tier 1: local model};
  \node[gate, below=4mm of ollama] (g1) {faithful?};
  \node[llm, below=4mm of g1] (gemini) {tier 2: cloud model};
  \node[gate, below=4mm of gemini] (g2) {faithful?};
  \node[det, below=4mm of g2] (rule) {tier 3: deterministic};
  \node[tier, fill=black!3, below=6mm of rule] (out) {summary returned};

  \draw[arr] (req)    -- (ollama);
  \draw[arr] (ollama) -- (g1);
  \draw[arr] (g1)     -- node[right, fail, xshift=1mm] {no, or tier failed} (gemini);
  \draw[arr] (gemini) -- (g2);
  \draw[arr] (g2)     -- node[right, fail, xshift=1mm] {no, or tier failed} (rule);
  \draw[arr] (rule)   -- (out);

  % Success paths bypass the remaining tiers.
  \draw[arr] (g1.west) -- ++(-9mm,0) node[above, note, xshift=2mm] {yes}
             |- (out.west);
  \draw[arr] (g2.west) -- ++(-5mm,0) node[above, note, xshift=2mm] {yes}
             |- (out.west);

  \node[note, right=4mm of rule, text width=24mm, text=green!35!black]
       {no network dependency; cannot fail};

\end{tikzpicture}
